\documentclass[12pt,a4paper]{article}
\usepackage[utf8]{inputenc}
\usepackage[spanish]{babel}
\usepackage{amsmath}
\usepackage{hyperref}
\usepackage{geometry}
\geometry{margin=2.5cm}

\title{\textbf{Sistema de Puntuación de Horarios (MatricuLAB v2)}}
\author{Documentación Técnica}
\date{\today}

\begin{document}

\maketitle

La asignación de puntos para determinar cuáles son los mejores horarios y qué porcentaje de ``perfección'' alcanzan se divide en dos procesos técnicos distintos: \textbf{Fase 1: Filtrado Heurístico (Web Worker)} y \textbf{Fase 2: Recálculo y Normalización Relativa (Frontend)}.

\section{Fase 1: Filtrado Heurístico Masivo (Web Worker)}

Cuando el usuario solicita generar horarios, el sistema evalúa millones de permutaciones mediante operaciones a nivel de bits (Bitwise). Durante este proceso en segundo plano, se genera un \textbf{Puntaje Bruto Absoluto (\texttt{rawScore})} para cada combinación válida. 

El único propósito de este puntaje bruto es alimentar una estructura de datos de cola de prioridad (\texttt{MinMaxHeap}) que retendrá únicamente los \textbf{50 mejores horarios en memoria RAM}, descartando el resto. Las reglas matemáticas aplicadas en esta fase son rígidas y priorizan la velocidad extrema de cálculo:

\begin{enumerate}
    \item \textbf{Menos Huecos (\texttt{min\_gaps}):}
    \begin{itemize}
        \item Se calculan todos los "slots" (bloques de 15 minutos) vacíos entre la primera y última clase de cada día.
        \item Penalización: Se resta el total de minutos de hueco de la semana ($\texttt{rawScore} \mathrel{-}= \text{totalGapMinutes}$).
    \end{itemize}
    
    \item \textbf{Menos Días (\texttt{min\_days}):}
    \begin{itemize}
        \item Se cuenta la cantidad de días en la semana con al menos una clase (\texttt{activeDaysCount}).
        \item Penalización: Se resta un peso muy alto por cada día físico requerido ($\texttt{rawScore} \mathrel{-}= \text{activeDaysCount} \times 500$).
    \end{itemize}

    \item \textbf{Evitar Días Huecos (\texttt{min\_day\_gaps}):}
    \begin{itemize}
        \item Se detectan los "días puente", es decir, días libres que están atrapados entre dos días de clases (por ejemplo, martes libre, pero lunes y miércoles con clases).
        \item Penalización: Castigo masivo por cada día puente detectado ($\texttt{rawScore} \mathrel{-}= \text{bridgeDays} \times 1000$).
    \end{itemize}

    \item \textbf{Preferir Mañana (\texttt{morning}):}
    \begin{itemize}
        \item Cada bloque de 15 minutos de clase que ocurra \textbf{estrictamente antes de las 12:00 PM} suma 1 a \texttt{morningSlots}.
        \item Bonificación: Se suman los minutos totales matutinos escalados ($\texttt{rawScore} \mathrel{+}= (\text{morningSlots} \times 15) / 5$).
    \end{itemize}

    \item \textbf{Preferir Tarde (\texttt{afternoon}):}
    \begin{itemize}
        \item Cada bloque de 15 minutos de clase que ocurra \textbf{a partir de las 12:00 PM en adelante} suma 1 a \texttt{afternoonSlots}.
        \item Bonificación: Se suman los minutos totales vespertinos escalados ($\texttt{rawScore} \mathrel{+}= (\text{afternoonSlots} \times 15) / 5$).
    \end{itemize}

    \item \textbf{Garantizar Almuerzo (\texttt{lunch}):}
    \begin{itemize}
        \item Se revisa si dentro de la ventana de almuerzo elegida existe un bloque libre consecutivo $\ge$ a la duración solicitada.
        \item Bonificación: Se saca la fracción de días que cumplen el requisito sobre el total de días activos y se premia fuertemente ($\texttt{rawScore} \mathrel{+}= \text{fracción} \times 1000$).
    \end{itemize}
\end{enumerate}

\section{Fase 2: Recálculo y Normalización Relativa (Frontend)}

Una vez que el Web Worker finaliza su ejecución masiva, envía el Top 50 de horarios al hilo principal (Frontend). En este punto, los puntajes brutos (\texttt{rawScore}) son descartados, ya que carecen de contexto humano intuitivo. El Frontend aplica su propia matemática refinada para calcular el \textbf{Porcentaje de Afinidad (0\% a 100\%)} final.

\subsection{2.1. Recálculo Elástico de Mañana y Tarde}
Para brindar mayor precisión y castigar severamente las clases que invadan el turno contrario, el frontend recalcula las métricas de mañana y tarde evaluando independientemente la hora de inicio de cada sesión (\texttt{startMinutes}):

\begin{itemize}
    \item \textbf{Métrica Mañana (\texttt{morningScore}):} \\
    Fórmula: $\max(0, 720 - \texttt{startMinutes})$. \\
    \textit{Significado:} Solo otorgan puntos las clases que inician antes de las 12:00 PM (720 min). Mientras más temprano inician (ej. 7:00 AM / 420 min), mayor es el puntaje ($720 - 420 = 300 \text{ pts}$). Cualquier clase desde el mediodía en adelante aporta $0$ puntos matutinos.
    
    \item \textbf{Métrica Tarde (\texttt{afternoonScore}):} \\
    Fórmula: $\max(0, \texttt{startMinutes} - 720)$. \\
    \textit{Significado:} Solo otorgan puntos las clases que inician a partir de las 12:00 PM. Mientras más tarde inician, mayor es el puntaje. Las clases matutinas aportan $0$ puntos vespertinos.
\end{itemize}

\subsection{2.2. Cálculo de Interpolación Lineal (Relativo al Top 50)}
Para determinar el porcentaje final de un horario, el sistema lo evalúa de forma \textbf{relativa frente a los otros 49 competidores} del lote:

\begin{enumerate}
    \item \textbf{Obtención de Límites:} Por cada filtro activado por el usuario, el algoritmo escanea el Top 50 en memoria y encuentra el \textbf{valor mínimo} y \textbf{máximo} absoluto de esa métrica.
    \item \textbf{Normalización por Métrica ($0$ a $1$):} Utilizando interpolación lineal, se ubica cada horario dentro de esos límites absolutos:
    \begin{equation}
        n = \frac{\text{valor\_horario} - \text{mínimo}}{\text{máximo} - \text{mínimo}}
    \end{equation}
    \begin{itemize}
        \item \textit{Filtros de Reducción (Huecos, Días):} Como se busca que el valor numérico sea lo menor posible, el cálculo se invierte ($1 - n$). El horario con menos huecos de todo el grupo recibe $1$ (100\%), y el que tiene más recibe $0$ (0\%).
        \item \textit{Filtros de Maximización (Mañana, Tarde):} El cálculo es directo. El horario con más horas acumuladas según la fórmula elástica recibe $1$ (100\%).
        \item \textit{Empate Global Matemático:} Si en todo el Top 50 el mínimo es idéntico al máximo (ej. todos los horarios tienen exactamente 2 días de clase), el divisor se anula. El sistema infiere perfección total en esa métrica y asigna $1$ (100\%) a todos.
    \end{itemize}
    \item \textbf{Porcentaje Final (\texttt{score}):} El sistema calcula el \textbf{promedio aritmético exacto} de las puntuaciones (de 0 a 1) obtenidas en todas las métricas objetivo, transformando el resultado en un porcentaje global final visual.
\end{enumerate}

\subsection{2.3. Jerarquía de Desempate (Tie-Breakers)}
Si tras la normalización, dos o más horarios empatan con un porcentaje final matemático exacto (ej. 85.00\%), el sistema interviene ordenándolos visualmente en la interfaz mediante reglas de desempate duro:
\begin{enumerate}
    \item Menor cantidad de minutos de hueco en total en la semana.
    \item Menor cantidad de días de asistencia requeridos.
    \item El horario cuya última sesión de clase termine más temprano a nivel semanal (\texttt{latestEndMinutes}).
\end{enumerate}

\end{document}
